\section{Capacidad de comprensión 3D de los modelos de video}

\subsection{Ventajas de los modelos generadores 2D}

En contraste marcado con los modelos generadores 3D, los modelos de generación de imágenes y videos en 2D pueden escalar eficazmente (scale) con los datos y los recursos computacionales: hay miles de millones de imágenes y videos en 2D disponibles en Internet para entrenarlos, y estos modelos generan píxeles 2D directamente, con sesgos inductivos limitados. A medida que crece el poder computacional, la calidad de los videos producidos por estos modelos mejora continuamente; por ejemplo, el modelo V3 ya puede generar videos con movimientos de cámara diversos, movimientos de escena y efectos físicos altamente realistas.

\subsection{¿Entienden los modelos de video el 3D?}

Una pregunta natural es: ¿son físicamente correctos los resultados generados por estos modelos de video?

El equipo de Henzler diseñó un experimento para investigar este problema:

\begin{enumerate}
    \item Se dan dos imágenes de escena, A y B, con áreas superpuestas muy reducidas
    \item Se utilizan estimadores de pose existentes (como DUSt3R) para predecir la cámara pose; debido a la poca superposición de las imágenes, los estimadores basados en correspondencias visuales funcionan mal
    \item Se interpolan fotogramas intermedios entre ambas imágenes utilizando diversos modelos generadores de video
    \item Los fotogramas intermedios generados se ingresan junto con las dos imágenes originales al estimador de pose
\end{enumerate}

\begin{importantbox}{Hallazgo clave: los modelos de video comprenden implícitamente el 3D}
Los resultados del experimento demuestran que los fotogramas intermedios generados por los modelos de video mejoran significativamente la precisión del estimador de pose. Esto significa que los modelos de video sí pueden modelar implícitamente el mundo físico: existen relaciones geométricas 3D razonables entre sus fotogramas, incluso si el modelo nunca fue entrenado explícitamente para comprender estructuras 3D.
\end{importantbox}

\subsection{Modelos de video interactivos}

Además, modelos de video interactivos similares a Genie 2 ya permiten a los usuarios navegar en tiempo real por un mundo 3D. Los usuarios pueden controlar la generación de nuevos fotogramas mediante entradas de teclado (moverse adelante, girar a la izquierda, girar a la derecha), explorando libremente diferentes perspectivas en una escena virtual. Lo más importante es que, cuando el usuario gira para mirar hacia atrás, la escena mantiene su \textbf{consistencia}; el aula aparece exactamente igual que antes.

\subsection{¿Por qué sigue siendo necesaria la modelación 3D explícita?}

Dado que los modelos de video puros en 2D ya son tan potentes, ¿por qué sigue siendo necesario modelar explícitamente el 3D?

\begin{warningbox}{Cuello de botella fundamental de los modelos de video: costo de inferencia}
Aunque los modelos de video pueden escalar con los datos y la potencia computacional, su costo de inferencia también aumenta proporcionalmente. Esto impide que estos modelos funcionen actualmente en dispositivos móviles como teléfonos inteligentes y tabletas. Por el contrario, las representaciones 3D explícitas (como NeRF o 3D Gaussian Splatting), una vez construidas, tienen un costo de renderizado extremadamente bajo y pueden ejecutarse en tiempo real en dispositivos móviles.
\end{warningbox}

\subsection{Combinación de dos piezas del rompecabezas}

Henzler señala que ya disponemos de dos piezas complementarias del rompecabezas:

\begin{itemize}
    \item \textbf{Modelos generadores de imágenes/videos}: producen imágenes 2D altamente realistas
    \item \textbf{Métodos de reconstrucción NeRF/3DGS}: destilan múltiples imágenes en representaciones 3D eficientes, permitiendo un renderizado rápido
\end{itemize}

Combinar ambas —utilizar modelos generadores para producir imágenes multivistas y luego destilar esas imágenes mediante métodos de reconstrucción 3D para obtener representaciones 3D eficientes— es precisamente el enfoque central del trabajo posterior CAT3D y B3D.

\subsection{Resumen de la sección}

Los modelos de video ya han demostrado una capacidad implícita de comprensión 3D, pero su costo de inferencia limita el despliegue en dispositivos móviles. Las representaciones 3D explícitas ofrecen la ventaja de un renderizado de bajo costo; por lo tanto, «generación 2D + destilación 3D» se ha convertido en la estrategia óptima para equilibrar calidad y eficiencia.
